

\section{Pengembangan Model Prediksi}

\subsection{Skenario Pengujian Anti-Overfitting}
Untuk memastikan model dapat digeneralisasi dengan baik pada data baru, diterapkan strategi \textit{Anti-Overfitting} sebagai berikut:
\begin{itemize}
    \item \textbf{Pembagian Data Stratifikasi:} Menggunakan \texttt{train\_test\_split} dengan \texttt{stratify=y} untuk menjaga proporsi kelas pelanggaran (7,47\%) tetap sama baik di data latih maupun data uji.
    \item \textbf{Pembatasan Kompleksitas Model:} Pada algoritma berbasis pohon (\textit{Tree-based}), kedalaman pohon (\texttt{max\_depth}) dibatasi (misal: \texttt{max\_depth=3} pada XGBoost dan \texttt{max\_depth=6} pada Random Forest) untuk mencegah model ``menghafal'' data latih.
\end{itemize}

\subsection{Seleksi Model (\textit{Model Selection})}
Penulis melakukan percobaan pelatihan menggunakan beberapa algoritma yang dipilih berdasarkan kemampuannya menangani data tidak seimbang:
\begin{enumerate}
    \item \textbf{Decision Tree:} Algoritma pembelajaran dasar yang memodelkan pengambilan keputusan menggunakan struktur seperti pohon. Model ini dipilih sebagai \textit{baseline} karena kemudahannya untuk diinterpretasi (\textit{interpretability}) melalui aturan logika "jika-maka" yang jelas.
    \item \textbf{Balanced Random Forest (BRF):} Varian dari algoritma \textit{Random Forest} yang dirancang khusus untuk menangani data tidak seimbang (\textit{imbalanced data}). Algoritma ini melakukan \textit{down-sampling} secara acak pada kelas mayoritas di setiap pembentukan pohon (\textit{bootstrap}) untuk menyeimbangkan distribusi kelas.
    \item \textbf{XGBoost (\textit{Extreme Gradient Boosting}):} Implementasi algoritma \textit{gradient boosting} yang sangat efisien, fleksibel, dan portabel. XGBoost dipilih karena kemampuannya dalam menangani data dalam skala besar dengan kecepatan eksekusi yang tinggi serta memiliki mekanisme regularisasi untuk mencegah \textit{overfitting}.
    \item \textbf{Support Vector Machine (SVM):} Algoritma yang bekerja dengan mencari \textit{hyperplane} pemisah optimal antar kelas. Dalam penelitian ini, digunakan kernel RBF (\textit{Radial Basis Function}) untuk menangkap pola hubungan non-linier yang kompleks antar fitur.
    \item \textbf{K-Nearest Neighbors (KNN):} Algoritma klasifikasi berbasis instans (\textit{instance-based learning}) yang sederhana. KNN mengklasifikasikan data baru berdasarkan mayoritas label dari $k$ tetangga terdekatnya menggunakan perhitungan jarak (seperti Euclidean).
\end{enumerate}


\subsection{Analisis Keunggulan XGBoost}
Berdasarkan Tabel \ref{tab:model_comparison}, \textbf{XGBoost} dipilih sebagai model terbaik dengan F1-Score tertinggi (\textbf{0.9003}) pada \textit{threshold} 0.48. Meskipun algoritma lain seperti Balanced Random Forest (BRF) dan EasyEnsemble juga mampu mencapai Recall sempurna (1.0), XGBoost unggul tipis dari segi Presisi (0.8187 dibandingkan 0.8142).

Keunggulan XGBoost dalam eksperimen ini didasari oleh beberapa faktor teknis:
\begin{enumerate}
    \item \textbf{Mekanisme Boosting yang Efisien:} XGBoost (\textit{Extreme Gradient Boosting}) bekerja dengan memperbaiki kesalahan prediksi dari model sebelumnya secara iteratif. Mekanisme ini terbukti efektif membedakan pola penyewa mencurigakan yang memiliki batas keputusan (\textit{decision boundary}) kompleks, sehingga meminimalkan \textit{false positive} lebih baik daripada algoritma \textit{bagging}.
    
    \item \textbf{Stabilitas pada Threshold Moderat:} XGBoost mencapai performa optimal pada \textit{threshold} 0.48, yang mendekati nilai \textit{default} 0.5. Hal ini menunjukkan kalibrasi probabilitas yang baik dan tidak memerlukan penyesuaian ekstrem seperti SVM (0.19) atau BRF (0.33), yang menjadikannya lebih andal untuk operasional.
    
    \item \textbf{Regularisasi Intrinsik:} Adanya parameter regularisasi bawaan (L1 dan L2) pada XGBoost membantu mencegah \textit{overfitting}, menjaga konsistensi performa tinggi tanpa mengorbankan generalisasi pada data uji.
\end{enumerate}

\subsection{Analisis Penentuan Ambang Batas Keputusan (\textit{Threshold Tuning})}
Untuk mendukung kebutuhan operasional, fungsi \texttt{find\_thresholds} digunakan guna mencari titik optimal. Pada konfigurasi XGBoost terpilih (Threshold 0.48), model mencapai keseimbangan terbaik:
\begin{itemize}
    \item \textbf{Recall 1.0000:} Model berhasil mendeteksi \textbf{seluruh} (100\%) penyewa mencurigakan dalam data uji, memastikan tidak ada ancaman keamanan yang terlewat (\textit{zero false negative}).
    \item \textbf{Precision 0.8187:} Tingkat akurasi prediksi positif mencapai 81,8\%, yang berarti tingkat notifikasi palsu (\textit{false alarm}) berada dalam batas wajar dan dapat diterima.
\end{itemize}

\subsection{Analisis Fitur Penting (\textit{Feature Importance})}
Analisis terhadap model XGBoost mengungkap bahwa fitur temporal adalah prediktor utama anomali. Fitur dengan kontribusi terbesar meliputi:
\begin{itemize}
    \item \textbf{\texttt{duration\_hours}:} Durasi sewa yang sangat singkat sering kali menjadi indikator kuat aktivitas transit yang tidak wajar.
    \item \textbf{\texttt{checkin\_hour}:} Waktu kedatangan di jam-jam tidak lazim (misalnya dini hari) berkorelasi tinggi dengan potensi pelanggaran.
    \item \textbf{\texttt{metode\_pembayaran}:} Pola pembayaran tertentu turut berkontribusi dalam membedakan profil risiko penyewa.
\end{itemize}




\subsection{Evaluasi dan Perbandingan Kinerja}
Evaluasi dilakukan menggunakan metrik \textit{F1-Score} dan \textit{Recall} pada kelas minoritas (label 1), mengingat sifat data yang sangat tidak seimbang. Berikut adalah ringkasan hasil evaluasi model terbaik:

\begin{table}[H]
\centering
\caption{Perbandingan Performa Model pada Kelas Minoritas (Mencurigakan)}
\label{tab:model_comparison}
    \begin{tabular}{|l|c|c|c|c|}
    \hline
    \textbf{Model} & \textbf{Mode/Threshold} & \textbf{Precision (1)} & \textbf{Recall (1)} & \textbf{F1-Score (1)} \\
    \hline
    Balanced RF & Balanced (0.33) & 0.8142 & 1.0000 & 0.8976 \\
    \hline
    \textbf{XGBoost} & \textbf{Balanced (0.48)} & \textbf{0.8187} & \textbf{1.0000} & \textbf{0.9003} \\
    \hline
    EasyEnsemble & Balanced (0.27) & 0.8142 & 1.0000 & 0.8976 \\
    \hline
    SVM (RBF) & Balanced (0.19) & 0.8156 & 0.9799 & 0.8902 \\
    \hline
    KNN & Balanced (0.61) & 0.7872 & 0.4966 & 0.6091 \\
    \hline
    \end{tabular}
\end{table}

